%!TeX root=csp-proof.tex
\section{Status and imports}\label{sec:status}

This section says, for every statement in this document, whether it is proved
here, imported, or not yet proved. It replaces the running commentary that a
partial exposition invites, and it is the only place where completeness claims
are made.

\subsection{The five verdicts}\label{sec:verdicts}

\begin{description}\tightlist
\item[\textsc{proved}] A complete proof is given here, or the statement is an
  immediate specialisation of one that is.
\item[\textsc{imported}] Proved elsewhere and used as a black box. Every such
  statement carries an exact citation in \S\ref{sec:imports-ledger}, with the
  hypotheses in the form the source states them.
\item[\textsc{outlined}] An argument is given at the level of its case
  structure, or with steps suppressed. \emph{This is not a proof} and the
  statement should not be relied on.
\item[\textsc{stated}] The statement appears and nothing is offered towards it.
  These are gaps in the logical chain, not imports.
\item[\textsc{open}] Confirmed to need an argument that is not available.
  There are none left in the table; what remains of that kind is
  Claim~\ref{cl:multiply-step} and Claim~\ref{cl:typing-anchor}, both
  written into proofs rather than left implicit.
  Hypothesis~\ref{hyp:maximal-mult} was the third; since the proof of
  Lemma~\ref{lem:parallelogram-crucial} no longer routes through
  Lemma~\ref{lem:maximal-mult}, it is no longer on any path to the main
  statements (Caveat~\ref{haz:typing-anchor}).
\end{description}

Nothing below is marked \textsc{proved} on the strength of a machine search.
Two claims here rest on finite verification --- the counterexamples of
Remark~\ref{rem:phi-twist} and of Caveat~\ref{haz:ba-center-implies} --- and
both are written so that a reader can check them in a few lines by hand.

One structural property \emph{is} checked mechanically: the proof-dependency
graph of this document --- an edge from each statement to every statement its
proof cites --- is acyclic, over $116$ statements, $80$ of which carry proofs,
and $190$ edges. That is a check on the arrangement, not on the mathematics,
and it has caught three real cycles: Caveats~\ref{haz:reverse-hom-split} and
\ref{haz:central-critical-split}, and a third that a first draft of
\pkey{lem:self-intersection-pc} introduced by borrowing four steps from the
proof of \pkey{lem:intersection-good} --- which consumes it. The borrowing was
only rhetorical, but the parser is right that a proof may not cite a statement
that cites it, and the steps are now written out.

\subsection{The table}\label{sec:status-table}

\begin{center}\footnotesize
\begin{longtable}{@{}p{0.33\textwidth}p{0.13\textwidth}p{0.46\textwidth}@{}}
\hline
\textbf{Statement} & \textbf{Verdict} & \textbf{Note} \\
\hline\endhead
\multicolumn{3}{@{}l}{\textit{\S\ref{sec:conventions}--\S\ref{sec:types}: vocabulary}}\\
\pkey{prop:p-divides}, \pkey{cor:Zp-in-Vn} & \textsc{proved} & \\
\pkey{lem:sg-terms} & \textsc{imported} & standard; \cite{Bergman} \\
\pkey{lem:rect-basics}, \pkey{prop:cover}, \pkey{lem:irred-quotient},
  \pkey{lem:linked-basics} & \textsc{proved} & \\
\pkey{lem:bridge-comp} & \textsc{imported} & \cite[Lem.~6.3]{ZhukJACM} \\
\pkey{lem:symmetrise} & \textsc{proved} & the collapse identity is proved by
  unfolding, so no irreducibility hypothesis is borrowed \\
\pkey{lem:central-ternary} & \textsc{imported} & \cite[Cor.~6.11.1]{ZhukStrong} \\
\pkey{lem:sfree-equiv}, \pkey{lem:linear-bridge} & \textsc{proved} & \\
\pkey{lem:linear-equiv}, \pkey{lem:bridge-reflexive},
  \pkey{lem:rectangularise} & \textsc{proved} & \\
\pkey{lem:no-cross-bridge}, \pkey{lem:bridge-to-pc} & \textsc{proved} &
  written out in \S\ref{sec:bridge-proofs}; the two normalisations they rested
  on are now \pkey{lem:rectangularise} and three lines,
  Caveat~\ref{haz:rectangularise} \\
\pkey{lem:pc-bridges-trivial} & \textsc{proved} & third case not written out \\
\pkey{lem:linear-on-top} & \textsc{proved} & \\
\pkey{lem:pc-on-top} & \textsc{stated} & consumed by nothing, and now
  carrying the BA-and-centre-free hypothesis its proof needs;
  Caveat~\ref{haz:on-top} \\
\hline
\multicolumn{3}{@{}l}{\textit{\S\ref{sec:instances}: instances}}\\
\pkey{prop:covering-basics} & \textsc{proved} & (g) is the imported
  Lemma~\ref{lem:expanded-consistency}; (e), (f) carry the renaming proviso of
  Caveat~\ref{haz:renaming} \\
\pkey{lem:conone-basics} & \textsc{proved} & transitivity needs only
  rectangularity; reflexivity on all of $A_{i}$ is where subdirectness is
  spent \\
\pkey{def:constraint-relation} & --- & the repeated-variable policy; every
  variable-named operation below is taken on $R^{C}$,
  Caveat~\ref{haz:diagonal} \\
\pkey{def:violation}, \pkey{lem:weakening-\brk terminates},
  \pkey{rem:get-crucial} & \textsc{proved} & the measure counts \emph{excluded}
  assignments, and cruciality quantifies over \emph{single} weaker constraints;
  Caveat~\ref{haz:weakening} \\
\hline
\multicolumn{3}{@{}l}{\textit{\S\ref{sec:properties}--\S\ref{sec:strongproofs}: strong subuniverses}}\\
\pkey{thm:brady-absorption} & \textsc{imported} & \cite[Thm.~3.11.1]{BradyNotes} \\
\pkey{lem:linked-implies-abs} & \textsc{proved} & a corollary of the above \\
\pkey{lem:pp-propagation} & \textsc{imported} & \cite[Lem.~2.9]{DecidingAbsorption};
  conclusion dotted, call sites audited, Caveat~\ref{haz:pp-dotted} \\
\pkey{lem:barto-kazda} & \textsc{imported} & \cite[Prop.~2.14]{DecidingAbsorption} \\
\pkey{lem:factor-type} & \textsc{imported} & $\TC$ only; \cite[Lem.~6.8]{ZhukStrong} \\
\pkey{lem:power-to-base} & \textsc{imported} & \cite[Lem.~6.24]{ZhukStrong} \\
\pkey{thm:central-\brk relation-source} & \textsc{imported} & \cite[Thm.~6.15]{ZhukStrong} \\
\pkey{lem:hobby-mckenzie} & \textsc{imported} & \cite{HobbyMcKenzie} \\
\pkey{lem:possible-\brk intersections} & \textsc{imported} & \cite[Lem.~6.25]{ZhukStrong} \\
\pkey{lem:absorbing-\brk equality} & \textsc{imported} & \cite[Lem.~7.2]{ZhukJACM} \\
\pkey{lem:ba-center-\brk implies} & \textsc{imported} &
  \cite[Cor.~6.1.2, 6.9.2]{ZhukStrong}; both hypotheses restored,
  Caveat~\ref{haz:ba-center-implies} \\
\pkey{lem:special-wnu} & \textsc{imported} & \cite[Lem.~4.7]{MarotiMcKenzie} \\
\pkey{lem:building-perfect} & \textsc{imported} & \cite[Cor.~8.17.1]{ZhukJACM} \\
\pkey{lem:common-binary-witness} & \textsc{proved} & the missing hypothesis of
  \pkey{lem:ba-center-implies} is exactly productivity of binary absorption;
  Caveat~\ref{haz:witness-moved} \\
\pkey{lem:reverse-hom}, \pkey{lem:bacon-left},
  \pkey{lem:no-central-critical}, \pkey{cor:no-central-critical},
  \pkey{lem:bridge-abelian}, \pkey{lem:ba-center-intersect},
  \pkey{lem:no-absorbing-diagonal}, \pkey{lem:absorb-linked-sub},
  \pkey{lem:no-unary-quotient}, \pkey{lem:central-relation},
  \pkey{lem:strong-nonempty}, \pkey{lem:multiset-shift},
  \pkey{lem:tot-sym-full}, \pkey{lem:tot-sym-irred}, \pkey{lem:full-fibre},
  \pkey{cor:intersection-good}, \pkey{lem:parallelogram-D},
  \pkey{lem:nice-bridge-abelian}, \pkey{lem:nontrivial-bridge},
  \pkey{lem:pc-inductive-step}, \pkey{lem:bridge-between-congruences},
  \pkey{lem:two-stable}, \pkey{lem:cut-or-nothing}, \pkey{lem:main-existence},
  \pkey{lem:ubiquity}, \pkey{lem:intersect-all},
  \pkey{cor:intersect-all-weak}, \pkey{lem:propagation},
  \pkey{cor:prop-quotient}, \pkey{thm:stable-intersection},
  \pkey{cor:stable-intersection}, \pkey{lem:multitype-chain},
  \pkey{lem:perfect-from-linked}, \pkey{lem:no-abs-in-linear},
  \pkey{lem:ba-central-meets}
  & \textsc{proved} & \\
\pkey{lem:self-intersection-pc} & \textsc{proved} & consumed by
  \pkey{lem:inter\brk section-good}, subsubcase 2A3;
  Caveat~\ref{haz:self-intersection} \\
\pkey{lem:intersection-good} & \textsc{proved} & both parts written out, with
  the five properness splits and the two full-fibre steps;
  Caveat~\ref{haz:intersection-draft} \\
\pkey{lem:anchored-intersection} & \textsc{proved} & the anchored form of
  stable intersection: leave-one-out is manufactured by first failure along
  the multi-type chain; closes configuration $(\gamma)$ in the proof of
  \pkey{lem:parallelogram-crucial} and the mixed-type half of the source's
  $\mathrm D19$; Caveat~\ref{haz:anchored} \\
\pkey{lem:preserve-linked-\brk anchored} & \textsc{proved} & the provable core
  of \pkey{lem:preserve-linked}: with the anchor, the full/cut dichotomy on
  the linking congruence closes both ways; consumed by nothing yet;
  Caveat~\ref{haz:preserve-anchor} \\
\pkey{lem:block-good-bridge}, \pkey{lem:center-pushed-in},
  \pkey{lem:leftlinked-stay-full}, \pkey{lem:factor-by-delta}
  & \textsc{outlined} & \\
\pkey{lem:multiply-all-linear} & \textsc{outlined} & Case~1 and the
  reduction of Case~2 are proved; Claim~\ref{cl:multiply-step}, which is the
  whole difficulty, is not. Caveat~\ref{haz:multiply-citation} \\
\pkey{cor:prop-from-factor}, \pkey{cor:prop-times-cong},
  \pkey{cor:prop-relations}, \pkey{lem:preserve-linked} & \textsc{stated} &
  the last two are now consumed by
  \pkey{lem:parallelogram-\brk crucial}'s proof \\
\pkey{lem:typed-above} & \textsc{stated} & the $t=1$ case is proved in
  Caveat~\ref{haz:typed-above}; consumed by the good branch of
  \pkey{lem:parallelogram-crucial} \\
\pkey{lem:maximal-mult} & \textsc{proved} & \emph{under the added
  Hypothesis~\ref{hyp:maximal-mult}}; consumed by nothing since
  \pkey{lem:parallelogram-crucial} acquired its own proof, so the hypothesis
  is no longer a debt of the chain;
  Caveats~\ref{haz:maximal-mult-hyp}, \ref{haz:maximal-mult} and
  \ref{haz:typing-anchor} \\
\hline
\multicolumn{3}{@{}l}{\textit{\S\ref{sec:main-statements}: the main statements}}\\
\pkey{thm:main-inductive}(2) & \textsc{proved} & modulo the common-term
  obligation of Caveat~\ref{haz:ba-center-implies} \\
\pkey{thm:main-inductive}(1) & \textsc{outlined} & structure only. The
  measure \emph{is} well founded, Caveat~\ref{haz:the-induction}(c) \\
\pkey{thm:reductions-safe} & \textsc{outlined} & the $\TBA$ and $\TC$ cases
  follow from \pkey{thm:main-inductive}(2); the $\TPC$ case is
  \pkey{thm:pc-safe}, which is stated only \\
\pkey{thm:pc-safe} & \textsc{stated} & \\
\pkey{lem:subuniverse-coset} & \textsc{proved} & \\
\pkey{thm:codim-one} & \textsc{outlined} & Steps 2--4 are now written out and
  the linear algebra is correct; what is outlined is Step~1, which rests on
  \pkey{lem:connected}, and two facts taken from the unread half.
  Caveat~\ref{haz:step2-linear-algebra} \\
\pkey{lem:expanded-\brk consistency} & \textsc{imported} & \cite[Lem.~6.1]{ZhukJACM} \\
\pkey{lem:tree-coverings} & \textsc{imported} & \cite[Lem.~5.6]{ZhukStrong} \\
\pkey{lem:crucial-irreducible}, \pkey{lem:bridge-from-relation} &
  \textsc{proved} & both were paraphrases before; the first had
  rectangularity as a conclusion where it is a hypothesis,
  Caveat~\ref{haz:crucial-irreducible} \\
\pkey{lem:connected},
  \pkey{lem:find-consistent},
  \pkey{lem:bridge-from-instance}, \pkey{cor:same-type} & \textsc{stated} &
  Caveat~\ref{haz:bridge-from-instance} \\
\pkey{lem:minimal-consistent} & \textsc{proved} & from
  Corollary~\ref{cor:prop-relations}(m1) and outright minimality;
  $1$-consistency of $\mathcal I$ restored as a hypothesis,
  Caveat~\ref{haz:minimal-consistent} \\
\pkey{lem:parallelogram-crucial} & \textsc{outlined} & the parallelogram
  property, irreducibility, the one-class fact, the good branch of the
  typing, and the exclusion of the $\TBA$/$\TC$ corner (via
  \pkey{lem:anchored-intersection}) are proved; what remains is
  Claim~\ref{cl:typing-anchor} and the stated \pkey{lem:typed-above}. Two hypotheses restored: subdirectness of
  $R$ and $\TS$-freeness of the $\mathbf D^{(1)}_{x_{i}}$,
  Caveats~\ref{haz:parallelogram-hyp} and \ref{haz:typing-anchor} \\
\hline
\multicolumn{3}{@{}l}{\textit{\S\ref{sec:algorithm}--\S\ref{sec:hardness}: the algorithm and the hard half}}\\
\pkey{def:solve-alg}, \pkey{lem:solve-\brk terminates} & \textsc{proved} &
  $\textsc{Solve}_{\mathrm{alg}}$ is now a function, by well-founded recursion
  on the measure of \pkey{def:solve-measure}; the three residual obligations
  are Caveat~\ref{haz:t1-names-solve} \\
\pkey{thm:t1} & \textsc{stated} & now has a subject; correctness itself is not
  proved here \\
\pkey{thm:t2}, \pkey{thm:h1} & \textsc{stated} & need a cost model \\
\pkey{thm:np-membership} & \textsc{stated} & needs the encoding and the
  machine; with \pkey{thm:h1} it supplies the word \emph{complete} \\
\pkey{lem:nae-interp} & \textsc{proved} & the interpretation written out from
  the two disjoint subalgebras \\
\pkey{thm:h0} & \textsc{proved} & from \cite[Cor.~1.5.11]{BradyNotes} and
  \pkey{lem:nae-interp} \\
\pkey{thm:galois} & \textsc{imported} & Geiger; Bodnarchuk et al. \\
\hline
\multicolumn{3}{@{}l}{\textit{\S\ref{sec:target}: the statement}}\\
\pkey{thm:dichotomy-core} & \textsc{stated} & the theorem this document is
  about; what it needs is everything above \\
\pkey{thm:core-reduction} & \textsc{imported} &
  \cite[Prop.~1.4.3, Thm.~1.4.7, Prop.~1.4.10, 1.4.11]{BradyNotes} \\
\pkey{lem:wnu-descends} & \textsc{proved} & \\
\pkey{thm:dichotomy} & \textsc{proved} & from \pkey{thm:dichotomy-core},
  \pkey{thm:core-\brk reduction} and \pkey{lem:wnu-descends};
  Caveat~\ref{haz:core} \\
\hline
\end{longtable}
\end{center}

\subsection{What this adds up to}\label{sec:status-summary}

The algebraic core of the strong-subuniverse theory
(\S\ref{sec:properties}--\S\ref{sec:strongproofs}) is the part in the best
shape: of its sixty-three statements, forty-one are proved here --- one of them,
\pkey{lem:maximal-mult}, only under Hypothesis~\ref{hyp:maximal-mult} ---
twelve are imported with exact citations, five are outlined, and five are
stated only. The two ends of the document, \S\ref{sec:target} and
\S\S\ref{sec:algorithm}--\ref{sec:hardness}, are now stated rather than
gestured at: the theorem is stated for a rigid core and derived for an
arbitrary $\Gamma$, $\textsc{Solve}_{\mathrm{alg}}$ is a function,
\pkey{thm:h0} is proved, and \emph{complete} has a membership statement behind
it. What is left there is the two complexity wrappers and the correctness of
the algorithm itself.
\S\ref{sec:main-statements} is in much worse shape: its spine ---
\pkey{lem:bridge-from-instance}, and with it most of
Lemma~\ref{lem:connected} --- is stated without proof, and part (1) of the main
induction is an outline whose measure is not yet well founded.
\pkey{lem:parallelogram-crucial} is the exception: its parallelogram half is
proved outright, and its typing half is proved up to
Claim~\ref{cl:typing-anchor} and Lemma~\ref{lem:typed-above}.

Three items block the rest, and they are the honest reading of the table.

\begin{enumerate}\tightlist
\item Claim~\ref{cl:typing-anchor} and Lemma~\ref{lem:typed-above}, the
  residue of the typing half of \pkey{lem:parallelogram-\brk crucial}. They
  replace Hypothesis~\ref{hyp:maximal-mult} in this list: the new proof does
  not cite \pkey{lem:maximal-mult}, so that hypothesis --- to which three
  routes failed, Caveat~\ref{haz:maximal-mult} --- is no longer on any path
  to the main statements. The claim cannot be settled at the level of
  subsets of a single algebra (the three-lines example of
  Caveat~\ref{haz:typing-anchor}), which is also why the failed routes
  failed.
\item The rewrite of \pkey{lem:multiply-\brk all-linear}, which is not yet a
  proof. With \pkey{lem:inter\brk section-good} and
  \pkey{lem:self-\brk intersection-pc} both written out, this is what is left
  of the intersection-and-factorisation layer.
\item Step~1 of \pkey{thm:codim-one} and the two facts its Steps~2--4 take
  from the unread half --- the linear algebra itself is now correct
  (Caveat~\ref{haz:step2-linear-algebra}); and, for the common binary
  absorption witness, the check that every producer of a type-$\TBA$ reduction
  produces a \emph{uniform} one, now that \pkey{thm:main-inductive}(2)
  hypothesises it (Caveats~\ref{haz:uniform-witness} and
  \ref{haz:witness-moved}).
\item[] \emph{Three entries stood here and no longer do.} The rectangularity
  normalisation is Lemma~\ref{lem:rectangularise}, and the second normalisation
  beside it is three lines. The induction measure
  across expanded coverings is not an open item --- no appeal is made at a
  covering, and Caveat~\ref{haz:the-induction}(c) enumerates the seven that
  are. The set-valued iteration behind $\Omega$ is
  Proposition~\ref{constr:omega}, proved once condition 3 is stated with
  single weakening steps.
\end{enumerate}

The dependency order among the three large statements is
\[
  \text{\pkey{thm:stable-intersection}}\;\longrightarrow\;
  \text{\pkey{lem:bridge-from-instance}}\;\longrightarrow\;
  \text{\pkey{thm:main-inductive}} .
\]
Stable intersection is the sole source of bridges; bridge-from-instance is the
only thing that turns them into instance-level structure; the main induction
consumes that. Nothing in the reverse direction is used.

\subsection{The import ledger}\label{sec:imports-ledger}

Every statement used and not proved here, with its source and the hypotheses in
the form the source states them. Where our statement differs from the printed
one --- because a hypothesis has been restored, or an alternative added --- the
difference is noted, and the reasons are in the companion audit document.

\begin{center}\footnotesize
\begin{longtable}{@{}p{0.28\textwidth}p{0.24\textwidth}p{0.40\textwidth}@{}}
\hline
\textbf{Here} & \textbf{Source} & \textbf{Hypotheses, and any difference} \\
\hline\endhead
\pkey{thm:brady-absorption} & \cite[Thm.~3.11.1]{BradyNotes} &
  $R\sdle\mathbf A\times\mathbf B$ subdirect and linked, $\mathbf A,\mathbf B$
  finite idempotent Taylor. Stated here in the source's three-alternative form;
  \pkey{lem:linked-implies-abs} is our weaker symmetric \emph{corollary} under
  properness, and is not Brady's theorem. \\
\pkey{lem:pp-propagation} & \cite[Lem.~2.9]{DecidingAbsorption},
  \cite[Lem.~6.1, Thm.~6.9]{ZhukStrong} &
  $R$ pp-defined by $\Phi$ in which $S$ occurs; $S'\subt{T}S$ with
  $T\in\{\TBA,\TC\}$ replaces \emph{every} occurrence. Note the conclusion is
  \emph{dotted}: replacing a relation by a proper strong subrelation can make
  the pp-defined result empty, so $R'\dsubte{T}R$ unless nonemptiness of $R'$
  is separately known. Every use below is in dotted form. \\
\pkey{lem:barto-kazda} & \cite[Prop.~2.14]{DecidingAbsorption},
  \cite[Lem.~3.2]{ZhukStrong} &
  $B\le\mathbf A$, $n\ge2$. Stated as printed. \\
\pkey{lem:factor-type} & \cite[Lem.~6.8]{ZhukStrong} &
  $B\subte{\TC}\mathbf A$, $\sigma$ a congruence on $\mathbf A$. Only the
  $\TC$ case is imported; $\TBA$ and $\TS$ are proved here
  (Caveat~\ref{haz:factor-type}). \\
\pkey{lem:power-to-base} & \cite[Lem.~6.24]{ZhukStrong} &
  $R$ a nontrivial strong subuniverse of $\mathbf A_{1}\times\dots\times
  \mathbf A_{n}$ of type $T\neq\TPC$. Already general in $T$; the case $T=\TS$
  needs no separate argument (Caveat~\ref{haz:power-to-base}). \\
\pkey{thm:central-relation-source} & \cite[Thm.~6.15]{ZhukStrong} &
  $R\sdle\mathbf A\times\mathbf B$, $\mathbf A,\mathbf B$ finite idempotent.
  Three alternatives, the third being a nontrivial projective subuniverse of
  $\mathbf B$; \pkey{lem:central-relation} eliminates it under
  Convention~\ref{conv:standing}, using also \cite[Lem.~3.4]{ZhukStrong}. \\
\pkey{lem:possible-intersections} & \cite[Lem.~6.25]{ZhukStrong} &
  $B\subt{T_{1}}\mathbf A$, $C\subt{T_{2}}\mathbf A$, $B\cap C=\varnothing$,
  $T_{1},T_{2}\in\{\TBA,\TC,\TS\}$. Stated as printed. \\
\pkey{lem:ba-center-implies} & \cite[Cor.~6.1.2, 6.9.2]{ZhukStrong} &
  $R\le\mathbf A_{1}\times\dots\times\mathbf A_{m}$ with
  $\proj_{1}(R)=A_{1}$, and $C_{i}\subte{\TBA(t)}\mathbf A_{i}$ for a
  \emph{single} binary $t$. Both hypotheses are as in
  \cite{ZhukStrong} and both are needed
  (Caveat~\ref{haz:ba-center-implies}). \\
\pkey{lem:hobby-mckenzie} & \cite{HobbyMcKenzie} &
  Abelian iff some congruence of the square has the diagonal as a block;
  Abelian iff affine, for a finite algebra with a WNU term. \\
\pkey{lem:absorbing-equality} & \cite[Lem.~7.2]{ZhukJACM} &
  $0_{\mathbf A}\subseteq\sigma\le\mathbf A^{2}$, $\omega\subt{\TBA}\sigma$
  \emph{proper}. The equality case needs separate treatment and gets it at the
  one call site. \\
\pkey{lem:bridge-comp} & \cite[Lem.~6.3]{ZhukJACM} &
  $\sigma_{0},\sigma_{1},\sigma_{2}$ \emph{all} irreducible. The hypothesis is
  easy to drop by accident; \pkey{lem:symmetrise} avoids the citation
  altogether. \\
\pkey{lem:building-perfect} & \cite[Cor.~8.17.1]{ZhukJACM} &
  $\sigma$ irreducible on $\mathbf A\in\Vn$, $\delta$ a bridge from $\sigma$ to
  $\sigma$ with $\bcol\delta=A^{2}$. Used only by
  \pkey{lem:perfect-from-linked}. \\
\pkey{lem:special-wnu} & \cite[Lem.~4.7]{MarotiMcKenzie} &
  $w$ an idempotent WNU of arity $n$ on a finite $A$. We use only
  ``$\Clo(w)$ contains a special idempotent WNU of some arity $N$'', not the
  printed bound. \\
\pkey{lem:expanded-consistency} & \cite[Lem.~6.1]{ZhukJACM} &
  an expanded covering of a cycle-consistent irreducible instance. \\
\pkey{lem:tree-coverings} & \cite[Lem.~5.6]{ZhukStrong} &
  a $1$-consistent instance and a reduction. \\
\pkey{lem:sg-terms} & \cite{Bergman} &
  standard; the ``fixed generator list'' phrasing is the one used
  (Caveat~\ref{haz:sg-fixed-list}). \\
\pkey{lem:central-ternary} & \cite[Cor.~6.11.1]{ZhukStrong} &
  a central subuniverse is ternary absorbing. \\
\pkey{thm:galois} & Geiger; Bodnarchuk--Kalu\v znin--Kotov--Romov &
  $A$ finite. See Caveat~\ref{haz:galois}. \\
Essential relations for $n\ge3$ & \cite[Lem.~6.11]{ZhukStrong} &
  $C_{i}\subte{\TC}\mathbf A_{i}$, $n\ge3$: no
  $(C_{1},\dots,C_{n})$-essential $R$. Used once, in
  \pkey{thm:stable-intersection}; its printed proof cites itself where it means
  the preceding lemma (Caveat~\ref{haz:d4}). \\
Cores and idempotent reducts & \cite{BradyNotes} &
  used by Caveat~\ref{haz:core} and \S\ref{sec:hardness}. \\
\hline
\end{longtable}
\end{center}

That is nineteen imported statements. The two that are real work are
\pkey{thm:brady-\brk absorption} and
\pkey{thm:central-\brk relation-source}; the rest are short or routine.
